% Created 2018-05-28 Mon 15:34
% Intended LaTeX compiler: pdflatex
\documentclass[11pt]{article}
\usepackage[utf8]{inputenc}
\usepackage[T1]{fontenc}
\usepackage{graphicx}
\usepackage{grffile}
\usepackage{longtable}
\usepackage{wrapfig}
\usepackage{rotating}
\usepackage[normalem]{ulem}
\usepackage{amsmath}
\usepackage{textcomp}
\usepackage{amssymb}
\usepackage{capt-of}
\usepackage{hyperref}
\usepackage[margin=0.80in]{geometry}
\usepackage[url=false, backend=biber, style=ieee]{biblatex}
\usepackage{float}
\usepackage[super,negative]{nth}
\usepackage[capitalise]{cleveref}
\usepackage{pst-node,transparent,ragged2e}
\addbibresource{/home/nuk3/.spacemacs.d/org-files/bibliography.bib}
\DeclareFieldFormat[inproceedings]{citetitle}{\textit{#1}}
\DeclareFieldFormat[inproceedings]{title}{\textit{#1}}
\DeclareFieldFormat[misc]{citetitle}{#1}
\DeclareFieldFormat[misc]{title}{#1}
\renewcommand*{\bibpagespunct}{%
\ifentrytype{inproceedings}
{\addspace}
{\addcomma\space}}
\AtEveryCitekey{\ifuseauthor{}{\clearname{author}}}
\AtEveryBibitem{\ifuseauthor{}{\clearname{author}}}
\SpecialCoor
\author{Nyameko Lisa}
\date{18 May 2018}
\title{IPSP055 - Essential Industrial Design Law\\\medskip
\large Assignment 03}
\hypersetup{
 pdfauthor={Nyameko Lisa},
 pdftitle={IPSP055 - Essential Industrial Design Law},
 pdfkeywords={},
 pdfsubject={},
 pdfcreator={Emacs 27.0.50 (Org mode 9.1.2)},
 pdflang={English}}
\begin{document}

\maketitle
\addvspace{110pt}
\centering{
\pnode(0.5\textwidth,-0.5\textheight){thisCenter}
\rput(thisCenter){%\transparent{0.25}
\includegraphics[width=2.7in]{/home/nuk3/course/llb/wipo-unisa/UNISACoatofArms.eps}}}

\justifying
\addvspace{110pt}
\section*{Declaration}
\label{sec:org9674b60}
I know that plagiarism is to use someone else’s work and pass it off as my own.\\
I know that plagiarism is wrong.\\
I confirm that this assignment is my own work.\\
I have acknowledged in the bibliography accompanying the assignment all the sources that I have used.\\
I have not directly copied without acknowledgement anything from the Internet or from any other source.\\
I have indicated every quotation and citation in a footnote or bracket linked to that quotation.\\
I have not allowed anyone else to copy my work and to pass it off as their own work.\\
I understand that if any unacknowledged copying whatsoever appears in my assignment I will receive zero per cent for the assignment.\\
I am aware of the UNISA policy on plagiarism and understand that disciplinary proceedings can be instituted against me by UNISA if I contravene this policy.\\
I indicate my understanding and acceptance of this declaration by
entering my name here under:
\begin{itemize}
\item Name: \textbf{Nyameko Lisa} (Student Number: \textbf{7874-909-3})
\end{itemize}

\subsection*{NOTE}
\label{sec:org3c268b4}
Please note that footnotes will be denoted as \footnote{This is a footnote} and will
appear at the bottom of the page.\\
References will be denoted by \cite{rsa93_designs_act} and will appear at the end of the document.
\newpage

\section{Will the \textbf{Vision 100 Motorrad} design qualify substantially for either functional or aesthetic design protection in South Africa?}
\label{sec:org1190d9f}
As per the definitions of [section 2(1)]\cite{rsa93_designs_act} and given that the
Republic is both a Paris Convention as well as a World Intellectual Property
Organisation (WIPO) Member State, [article 15(1)]\cite{wto17_trips}, [article
1(2)]\cite{wipo83_paris_conve_protect_ip}, \cite{wipo25_hague_agreement}, South African statutory design
legislation is in compliance with International Conventions and legal instruments.

\subsection{Requirements for design protection}
\label{sec:org7cf2b7b}

Holger Hampf, acting on behalf of BMW, who will henceforth be referred to as the
\textbf{``proprietor''} of the Vision 100 Motorrad design which will in turn be
referred to as the \textbf{``article''}, as per the provisions of [section
1(1)]\cite{rsa93_designs_act}. BMW and Holger's design and customer experience team,
should be advised that the legislation makes provision for both \uline{aesthetic
designs} (Part A) as well as \uline{functional designs} (Part F) registrations, as per
the provisions of [section 14(1)(a-b)]\cite{rsa93_designs_act}, for the protection
of the physical form of articles of manufacture intended to be multiplied
by industrial processes, [section 14(4)]\cite{rsa93_designs_act}.

\subsubsection{Aesthetic designs}
\label{sec:orga47b203}

As per the provisions of [section 1(1)]\cite{rsa93_designs_act}, which stipulates
that \textit{``...any design applied to any article, whether for the
pattern or the shape or the configuration or the ornamentation thereof, or for
any two or more of these purposes, and by whatever means is applied, having
features which appeal to and are judged solely by the eye, irrespective of the
aesthetic quality thereof...''}, as was stipulated by Judge Corbett in
\citetitle{corbett84_homecraft_steel_v_sm_hare} \cite{corbett84_homecraft_steel_v_sm_hare}:
\begin{quote}
\textit{To qualify as a design, an article must have some special feature
which appeals to the eye and distinguishes it from others of its class.''}
\end{quote}
Based on the images supplied in \textbf{Figures 1 \& 3}, the \textbf{proprietor} would be
advised to pursue aesthetic design registrations for a number of aesthetic design
elements of their \textbf{article}.

\subsubsection{Functional designs}
\label{sec:org87847c9}

Again as per the provisions of [section 1(1)]\cite{rsa93_designs_act}, the
definition of a functional design is \textit{``...any design applied to any article,
either for the pattern or the shape or the configuration thereof, or for any two
or more of these purposes, and by whatever means it is applied, having features
which are necessitated by the function which the article to which the design is
applied, is to perform, and includes an integrated circuit topography, a mask
work and a series of mask works...''}.\\

It is recommended that the \textbf{proprietor} lodge both functional and aesthetic
design registration applications in respect of their \textbf{article}, as per the
provisions of [section 14(1)(b)(ii)]\cite{rsa93_designs_act}. From the description
of the Vision 100 Motorrad, there are a number of purely functional design features, methods and
principles of it's construction [section 14(4)]\cite{rsa93_designs_act}, \emph{``that
are not commonplace in the art in question''}, requiring purely functional
design protection. Such as:
\begin{itemize}
\item Advanced technology to enhance safety and the thrill of riding, event allowing
for safe helmet-less riding,
\item Expanding block housing the bike's electric motor,
\item Reduction of mechanical elements, brake cables and buttons,
\item Triangular frame and it's bendable `Flexframe' material.
\end{itemize}

\subsubsection{Originality}
\label{sec:orgca78f25}

As per the provisions of [section 14(1)(a)]\cite{rsa93_designs_act}, should Holger
and BMW be able to demonstrate that the design features of their Vision 100
Motorrad are both \textbf{new} and \textbf{original}, they may in the prescribed manner and on
payment of the prescribed fee, apply for aesthetic and functional design
registrations.\\

Moreover the Vision 100 Motorrad must be different from or not form part of the
\textbf{state of the art} within a grace period of six months from the date of first
disclosure, or release date [section 14(2)(b)]\cite{rsa93_designs_act}. Where
state of the art comprises of all matter which has been made available to the
public, within or outside the Republic, by written means, usage or any other
form [section 14(3)(a)]\cite{rsa93_designs_act}, including all matter contained in
pending registration applications [section 14(3)(b)]\cite{rsa93_designs_act}.\\

Courts of the Republic have interpreted the meaning of \emph{`original'} to coincide
with that definition outlined in \cite{rsa78_copyrightact}, and which was further
elaborated on by Judge Van Heerden in
\citetitle{vheerden71_xactics_v_tailored_containers}
\cite{vheerden71_xactics_v_tailored_containers}, where he explained that:
\begin{quote}
\textit{``I have accordingly come to the conclusion that applicant's registered design
68/346 was not new or original by reason of the fact that a design not substantially
different therefrom had already been described in a printed publication before the
date of application for registration.''}
\end{quote}

Lastly in respect of the prior art depicted in \textbf{Figure 2}, it could be argued
that the commonplace in the art elements and aspects demonstrated therein,
describe features that are necessarily determined by the function of the article
in question and are \textbf{not} registrable as designs in and of themselves. In
relation to case law heard within the Republic, this was succinctly articulated by
Judge Advocate Corbett in \citetitle{corbett84_homecraft_steel_v_sm_hare}
\cite{corbett84_homecraft_steel_v_sm_hare},

\begin{quote}
\textit{``In terms of the definition of F "design" in s 1 (1) of the
Act, any feature which is dictated solely by the function which the article is
intended to perform, is excluded from protection under the Act. The top and
bottom flanges and the return lip on the bottom flange are purely functional in
that they provide the requisite rigidity and facilitate handling of the fascias.
The fluting on the face G of the design serves the functional purpose of
providing additional rigidity. In regard to the question of when functional
features are excluded from protection, see Amp Incorporated v Unilux (Pty) Ltd
[1972] RPC 103; Swisstool Manufacturing Co v Omega Africa Plastics (supra ). See
also the commentary on "method or principle of construction" in Russel-Clarke
(supra at 27 - 28) and Kestos Ltd v Kempat Ltd \& H Kemp [1936] 53 RPC 139 at
151.''}
\end{quote}

Irrespective of the fact that the Vission 100 incorporates registrable as well
as unregistrable features, it is the registratble features that the BMW would be
encouraged to pursue.

\section{Are there any other forms of intellectual property protection available to BMW for the Vision 100 Motorrad?}
\label{sec:orga652629}

Yes indeed a number of alternative forms of intellectual property protection
exist, namely \textbf{Patent}, \textbf{Trademark} and \textbf{Copyright} protection as discussed below.\\

As per the definition specified in [Article 27(1)]\cite{wto17_trips} and [section
2]\cite{rsa78_patents_act}, BMW's Vision 100 Motorrad need satisfy the provisions
for patents to be granted under [section 25]\cite{rsa78_patents_act}, wherein
subsection (1) of that section of the Act stipulates that a patent may be
granted for a \textbf{new} invention, involving an \textbf{inventive step} and has application
to trade or industry.\\

The content of the concept ``patentable subject-matter' is usually established
in the negative sense by statue, i.e. inventions explicitly `excluded' from
qualification for the purposes of patent protection. The HUD\footnote{Heads Up
Display.} display and operating concept, the expanding block house (and it's
controlling circuitry and semiconductor chip), the triangular frame and
`Flexframe' material, to name but a few, neither fall into any of the categories
listed under [Article 27(2) and Article 27(3)]\cite{wto17_trips}, nor
those articulated in [subsections (2) and (3) of section 25]\cite{rsa78_patents_act}. BMW
would be strongly advise to seek patent protection for various aspects of the
Vision 100 Motorrad.\\

Given that the Republic is a Paris Convention Member State, [article
15(1)]\cite{wto17_trips}, [article 1(2)]\cite{wipo83_paris_conve_protect_ip}, and
per the definitions of [section 2(1)]\cite{rsa93_tm_act}, it follows that \emph{`Vision
100 Motorrad'} as well as \emph{`Flexframe'} will be understood to be \textbf{marks}, in
that they are either names or signs capable of being represented graphically.
Moreover they constitute \textbf{trade marks}, as they have been proposed by BMW, in
relation to distinguishing their trade or use from other goods or services
connected in the course of trade regarding \emph{`class 12: motors and engines for
land vehicles'} and \emph{`class 6: common metals and their alloys'} respectively,
\cite{wipo57_ncl}. Therefore as per [section 9(1)]\cite{rsa93_tm_act}, BMW would be
advised to pursue trademark registration of the \textbf{Vision 100 Motorrad} and
\textbf{Flexframe} marks, subject to the limitations and provisions of [paragraphs (1),
(4), (5), (8), (9), (10) or (11) of section 10]\cite{rsa93_tm_act}.\\

Lastly as per the provisions of
\cite{rsa78_copyrightact,wto17_trips,wipo96_copyright_treaty,wipo86_berne},
copyright would subsist in any promotional material that BMW prepares in
relation to their Vision 100 Motorrad.

\section{State whether the actions of Carsmart International amount to infringement of the South African design for the Vision 100 Motorrad lights, to which Carsmart counterclaims.}
\label{sec:orgf82121c}

As per the provisions of [section 20(1)]\cite{rsa93_designs_act}, the effect of
registration of the design for replacement lights for the Vision 100 Motorrad in
\textbf{Figure 3} shall grant the registered proprietor in the Republic, for the
duration of registration [section 22(1)]\cite{rsa93_designs_act}, the exclusive
right to exclude others from \textbf{using}, \textbf{importing} or \textbf{disposing} \emph{(or
distributing)} of any article included in the class in which their design is
registered and embodying the registered design or a design not substantially
different from the registered design.

\subsection{Infringement}
\label{sec:orga70327c}

As per the provisions of [section 35(1)]\cite{rsa93_designs_act}, the registered
proprietor BMW, of the registered design in \textbf{Figure 3} may institute proceedings
for infringement against Carsmart International, for the \textbf{unauthorised
importation, use and disposal} of their registered design. Which by its very
nature \emph{`as a replacement part'}, must necessarily bear resemblance to, and not be
substantially different from BMW's registered design as an article of the same
class.

\subsection{Visual Judgement of Similarity}
\label{sec:orgfb09855}

As per \citetitle{corbett84_homecraft_steel_v_sm_hare}
\cite{corbett84_homecraft_steel_v_sm_hare}, the test for infringement is a purely
visual judgement according to the observation of the likely customer where Judge
Corbett referred to Judge Eloff in preceding case law,
\begin{quote}
\textit{``to hold that while the eye is to be the eye of the Court, the
Court should view the design through the spectacles of the customer''}
\end{quote}
wherein the designs are viewed and considered both independently as well as
side by side,
\begin{quote}
\textit{``an article must have some special feature which appeals to the
eye and distinguishes it from other of its class... This must be decided by the
eye; and in applying this visual test the Court must bear in mind the reaction
of the likely customer. The two designs must be viewed and compared side by side
and also separately.''}
\end{quote}

\subsection{Novelty and Originality}
\label{sec:orgaeda47e}

When analysing the possibility of infringement, the state of the art and the
degree of novelty and originality achieved by the registered design is also an
important factor, in reference to preceding case law, Judge Corbett stipulates,
\cite{corbett84_homecraft_steel_v_sm_hare} that,
\begin{quote}
\textit{``If only small differences seperate the registered design from
what has gone before, then eqaully small differences between the alleged
infringement and the regsitered design will be held to be sufficient to avoid
infringement...''}
\end{quote}
Given the very nature of the \textbf{infringing article}, in that it is a replacement
part, there must necessarily exist insignificant differences between the two
designs, in \textbf{shape} and/or \textbf{configuration}, and hence there is a high level of
risk that a member of the public or a likely customer will confuse the two
designs.

\subsection{Prescribed cases applicable to infringement and relief}
\label{sec:orgf4e9080}

Where in substantiating proceedings for design infringement applicants have in
the past relied on alleged infringement of their registered designs as well as
on the delict of passing off, \citetitle{corbett84_homecraft_steel_v_sm_hare}
\cite{corbett84_homecraft_steel_v_sm_hare},
\citetitle{vdwesthuizen02_slavepak_v_buddy} \cite{vdwesthuizen02_slavepak_v_buddy},
\citetitle{duplessis02_koninklijke_phillips_v_kenwood}
\cite{duplessis02_koninklijke_phillips_v_kenwood}. As per the provisions of
[section 35(3)(a-d respectively)]\cite{rsa93_designs_act}, the plaintiff BMW,
would likely seek relief in the form of:
\begin{itemize}
\item \textbf{an interdict} instructing Carsmart to cease importation, use and disposal of
the infringing articles,
\item \textbf{surrender} of the infringing replacement lights,
\item \textbf{damages} or \textbf{in lieu of damages, an amount calculated on the basis of a
reasonable licensee's royalty}, specifically with regards to unfair
competition resulting from \emph{passing of} as per [article 16(1)]\cite{wto17_trips}, [article
10\(^{\text{bis}}\)(3)(1)]\cite{wipo83_paris_conve_protect_ip} and [article
2(1)]\cite{wipo96_model_provi_unfair_comp};
and \emph{dilution of their goodwill and / or reputation} as per [article
3(2)]\cite{wipo96_model_provi_unfair_comp}.
\end{itemize}

\subsection{Counterclaim and revocation of design registration}
\label{sec:org58f26a5}

Given the above considerations it is clear that the actions of Carsmart
International amount to infringement, however as per the provisions of [section
35(5)]\cite{rsa93_designs_act}, Carsmart may counterclaim for the revocation of
the registration of the design. Furthermore Carsmart International may at any
time, in the prescribed manner apply to the court for the revocation of BMW's
design registration for the Vision 100 Motorrad lights, as per [section
31(1)(c)]\cite{rsa93_designs_act}. In both cases it would be argued that BMW's
aesthetic design for the lights is ineligible for registration in that it
constitutes a \textbf{spare part for a vehicle}, as per [section
14(6)]\cite{rsa93_designs_act}.\\

With this in mind BMW would be advised to voluntarily withdraw the registered
design for the Vision 100 Motorrad lights in question as per [section
34]\cite{rsa93_designs_act}.

\section{Is it possible to register a design right for the semiconductor chip that regulates the control of the `Flexframe' of the Vision 100 Motorrad?}
\label{sec:orgdcfe88f}

As per the provisions of [section 1]\cite{rsa93_designs_act}, [17 U.S.C.
901(a)(1,2)]\cite{usa47_title17_ch9} and [Article 2]\cite{wipo89_washington_treaty_ip_ic}
the semiconductor chip will be understood to be an \textbf{integrated circuit
(topography)} or \textbf{layout-design (topography)} or \textbf{mask work} or \textbf{semiconductor
chip product} constituting an article or pattern, shape or configuration of the
three-dimensional disposition, in final or intermediate form containing
electrical, electromagnetic or optical elements and circuitry, capable of
performing an opto-electromagnetic function in a semiconductor material.\\

As per the provisions of [17 U.S.C. 902(b)]\cite{usa47_title17_ch9} and [Article
3(2)]\cite{wipo89_washington_treaty_ip_ic} , in order fro BMW's semiconductor chip
to enjoy layout-design (topography) protection it must satisfy the \emph{requirement
of originality}, in that it\footnote{Or a combination of elements and
interconnections taken as a whole.} must be the result of Holger's and BMW's own
intellectual effort and not be commonplace among designers and manufactures of
integrated circuits at the time of its creation. Where as per [17 U.S.C.
904(a)]\cite{usa47_title17_ch9} and [Article
7(2)]\cite{wipo89_washington_treaty_ip_ic}, BMW may commence protecting their
semiconductor chip either upon first commercial exploitation or once
registration has been applied for.\\

In contrast to the International precedent discussed above, the Republic instead
introduces specific protection for integrated circuits, defined and protected as
functional designs as per [section 1]\cite{rsa93_designs_act}. Therefor it follows
that BMW's semiconductor chip will be registrable as a functional design if
Holger and BMW can demonstrate that it is \textbf{new} and \textbf{not commonplace in the art}
and \textbf{multiplied by and industrial process}, as per [sections 14(1)(b)(i-ii) and
24(4)]\cite{rsa93_designs_act} respectively.\\

With the above considerations, BMW would be advised to pursue intellectual
property protection of this semiconductor chip in Union Member States and the
United States, and to concurrently pursue \textbf{functional design protection} within
the Republic.

\printbibliography
\end{document}